\documentclass[10pt,a4paper]{amsart}
\usepackage[margin=19mm]{geometry}
\usepackage{amsmath,amssymb,mathtools}
\usepackage[scr=rsfs,cal=euler]{mathalfa}
\usepackage{xfrac}
\usepackage[unicode,hidelinks,bookmarks=false]{hyperref}
\newcommand{\ie}{i.e.\ }
\newcommand{\cf}{cf.\ }
\newcommand{\resp}{respectively\ }
\newcommand{\Subsection}{\subsection}
\providecommand{\nobreakdash}{\nobreak\hbox{-}}
\setlength{\emergencystretch}{2em}
\multlinegap0pt
\usepackage{amssymb}
\usepackage{xcolor}
\usepackage{enumerate}
\newtheorem{proposition}{Proposition}
\newtheorem{lemma}{Lemma}
\title{Generic Strong Well-Posedness for McKean--Vlasov SDEs with Jumps in a Krylov-Stable Coefficient Space}
\author{Mingbo Zhang}
\date{Source: arXiv:2609.06132v1 (CC BY 4.0)}
\begin{document}
\maketitle

\noindent\emph{Source and licence.} Generic Strong Well-Posedness for McKean--Vlasov SDEs with Jumps in a Krylov-Stable Coefficient Space. Version arXiv:2609.06132v1 (CC BY 4.0), distributed by arXiv under the Creative Commons Attribution 4.0 International licence, http://creativecommons.org/licenses/by/4.0/, as recorded in the arXiv licence metadata for this exact version and shown on the abstract page https://arxiv.org/abs/2609.06132v1. Full licence notice: \texttt{licenses/arxiv-2609.06132-CC-BY-4.0.txt}.\par\smallskip

\noindent\textit{Editorial scope.} Counted unit: the article's proposition thm:local-weighted-krylov (source0.tex lines 1244--1263). The article's complete original proof of it is delivered with it (source0.tex lines 1265--1971, one \texttt{proof} environment of 707 lines). No mathematics is written, repaired or re-derived: the statement and the proof are the article's own lines, in the article's own notation, and no retained line is substituted or re-flowed.  Retained uncounted as the source-local context the proof needs: lines 1233--1236; lines 7310--7371.  Nothing was omitted from any retained span, so the package carries no omission record.  No retained line contains a citation, so the wrapper carries no bibliography and no bibliography entry.  No retained line contains a figure environment, an image-inclusion command or drawing code, so no image is delivered and the package declares no asset dependency.  Editorial formatting only, in addition: an AMS \texttt{amsart} preamble that restates the article's own declarations and macros used by the retained lines, namely the environments \texttt{proposition}, \texttt{lemma} on separate counters, together with the standard packages those lines need.  The retained environments print in this document as Proposition 1, Lemma 1 in order of appearance, whereas the article prints the same items with the numbering of its own full document.  The complete native source of the article as distributed in the arXiv e-print for this exact version is bundled unchanged as \texttt{sources/209493/source0.tex}.
	\begin{equation}
		\label{eq:ito-levy-uniform-bounds}
		|b_s|\le K_b,\qquad\|\sigma_s\|_{\mathrm{HS}}\le K_\sigma,\qquad	\int_U|g_s(z)|^2\,\nu(dz)\le K_g^2.
	\end{equation}

	\begin{lemma}
		\label{lem:global-aleksandrov-regularization}
		Let $T,R>0$, and let
		\[
		f\in C_c^\infty((0,T)\times B_R),
		\qquad f\geq0.
		\]
		Let $\widetilde f$ denote the zero extension of $f$ to
		$\mathbb R\times\mathbb R^d$. Then there exist constants
		$\kappa_d>0$ and $C_{d,T,R}>0$, a sequence
		$\delta_n\downarrow0$, and functions
		\[
		u^n\in C^\infty(\mathbb R\times\mathbb R^d),
		\qquad
		f^n\in C_c^\infty(\mathbb R\times\mathbb R^d),
		\qquad
		f^n\geq0,
		\]
		such that the following properties hold:
		\begin{enumerate}
			\item[\textnormal{(i)}]
			for every $t\in\mathbb R$, the function
			$x\mapsto u^n(t,x)$ is convex on $\mathbb R^d$;
			
			\item[\textnormal{(ii)}]
			$f^n\to\widetilde f$ uniformly on
			$\mathbb R\times\mathbb R^d$;
			
			\item[\textnormal{(iii)}]
			for every symmetric non-negative definite matrix
			$A\in\mathbb R^{d\times d}$,
			\begin{equation}
				\label{eq:appendix-smoothed-aleksandrov-inequality}
				\partial_tu^n(t,x)
				+
				\operatorname{tr}\!\left(A D^2u^n(t,x)\right)
				\geq
				\kappa_d(\det A)^{1/(d+1)}f^n(t,x)
			\end{equation}
			for all $(t,x)\in[0,T]\times B_R$;
			
			\item[\textnormal{(iv)}]
			uniformly in $n$,
			\begin{equation}
				\label{eq:appendix-global-growth-bound}
				|u^n(t,x)|
				\leq
				C_{d,T,R}
				\|f\|_{L^{d+1}((0,T)\times B_R)}
				(1+|x|),
			\end{equation}
			and
			\begin{equation}
				\label{eq:appendix-global-gradient-bound}
				|\nabla u^n(t,x)|
				\leq
				C_{d,T,R}
				\|f\|_{L^{d+1}((0,T)\times B_R)}
			\end{equation}
			for all $(t,x)\in\mathbb R\times\mathbb R^d$.
		\end{enumerate}
	\end{lemma}

	\begin{proposition}
		\label{thm:local-weighted-krylov}
		Assume \eqref{eq:ito-levy-uniform-bounds}. Let $\gamma$ and $\tau$ be	stopping times satisfying
		\[
		0\le\gamma\le\tau\le\tau_R(Y).
		\]
		Then there exists a constant
		\[
		C_{\mathrm K}=C_{\mathrm K}(d,T,R,K_b,K_\sigma,K_g)
		\]
		such that, for every non-negative Borel function
		$f\in L^{d+1}(Q_{T,R})$,
		\begin{equation}
			\label{eq:weighted-conditional-krylov}
			\mathbb E\left[\left.\int_\gamma^\tau(\det A_s)^{1/(d+1)}f(s,Y_s)\,ds\,\right|\mathcal F_\gamma\right]
			\le
			C_{\mathrm K}\|f\|_{L^{d+1}(Q_{T,R})}
		\end{equation}
		almost surely. On the event $\{\gamma=\tau\}$ the left-hand side is understood to be zero.
	\end{proposition}

\begin{proof}
	We first prove the estimate for smooth functions and then pass to
	arbitrary \(L^{d+1}\)-functions by identifying the stopped occupation
	measure.
	
	
	\textbf{Step 1: smooth non-negative functions and the jump operator.}
	
	Let
	\[
	f\in C_c^\infty((0,T)\times B_R),
	\qquad
	f\geq0.
	\]
	By Lemma~\ref{lem:global-aleksandrov-regularization}, there exist
	\(u^n\in C^{1,2}([0,T]\times\mathbb R^d)\) and non-negative smooth
	functions \(f^n\) such that
	\begin{equation}
		\label{eq:regularized-f-convergence}
		\|f^n-f\|_{L^\infty(\mathbb R\times\mathbb R^d)}
		\longrightarrow0,
	\end{equation}
	\(x\mapsto u^n(t,x)\) is convex, and, for every symmetric
	non-negative definite matrix \(A\),
	\begin{equation}
		\label{eq:regularized-aleksandrov-inequality}
		\partial_tu^n(t,x)
		+\operatorname{tr}\!\left(A D^2u^n(t,x)\right)
		\geq
		\kappa_d(\det A)^{1/(d+1)}f^n(t,x)
	\end{equation}
	on \([0,T]\times B_R\), here \(D^2u\) denotes the Hessian with respect to the spatial variable. Moreover, with
	\[
	L_f
	:=
	C_{d,T,R}\|f\|_{L^{d+1}(Q_{T,R})},
	\]
	one has, uniformly in \(n\),
	\begin{equation}
		\label{eq:regularized-growth-gradient-bounds}
		|u^n(t,x)|
		\leq
		L_f(1+|x|),
		\qquad
		|\nabla u^n(t,x)|
		\leq
		L_f.
	\end{equation}
	
	For fixed \(n\), set
	\begin{equation}
		\label{eq:jump-operator}
		\mathcal J_su^n(t,x)
		:=
		\int_U r_n(t,x,g_s(z))\,\nu(dz),
	\end{equation}
	where
	\[
	r_n(t,x,h)
	:=
	u^n(t,x+h)-u^n(t,x)-\nabla u^n(t,x)\cdot h.
	\]
	Convexity gives
	\begin{equation}
		\label{eq:jump-operator-positive}
		r_n(t,x,h)\geq0,
		\qquad
		\mathcal J_su^n(t,x)\geq0.
	\end{equation}
	
	We next verify that the integral in
	\eqref{eq:jump-operator} is finite along the trajectory before the
	exit time. Put
	\[
	H_{n,R}
	:=
	\sup_{(t,x)\in[0,T]\times\overline B_{R+1}}
	\|D^2u^n(t,x)\|
	<\infty.
	\]
	If \(x\in\overline B_R\) and \(|h|\leq1\), Taylor's formula and
	convexity imply
	\begin{equation}
		\label{eq:small-jump-remainder}
		0
		\leq
		r_n(t,x,h)
		\leq
		\frac12 H_{n,R}|h|^2.
	\end{equation}
	If \(|h|>1\), the global gradient bound in
	\eqref{eq:regularized-growth-gradient-bounds} gives
	\begin{equation}
		\label{eq:large-jump-remainder}
		0
		\leq
		r_n(t,x,h)
		\leq
		2L_f|h|
		\leq
		2L_f|h|^2.
	\end{equation}
	Consequently, for \(x\in\overline B_R\),
	\begin{equation}
		\label{eq:jump-operator-integrability}
		0
		\leq
		\mathcal J_su^n(t,x)
		\leq
		\left(
		\frac12 H_{n,R}+2L_f
		\right)
		\int_U|g_s(z)|^2\,\nu(dz)
		<\infty.
	\end{equation}
	This small-jump/large-jump decomposition is the point at which the
	second-moment assumption on \(g\) is used to control the compensator
	term.
	
	
	\textbf{Step 2: the stochastic integrals are true martingales.}
	
	Define
	\[
	M_t^{W,n}
	:=
	\int_0^t
	\nabla u^n(s,Y_{s-})\sigma_s\,dW_s
	\]
	and
	\[
	M_t^{N,n}
	:=
	\int_0^t\int_U
	\bigl[
	u^n(s,Y_{s-}+g_s(z))
	-u^n(s,Y_{s-})
	\bigr]
	\,\widetilde N(ds,dz).
	\]
	By \eqref{eq:regularized-growth-gradient-bounds} and
	\eqref{eq:ito-levy-uniform-bounds},
	\begin{equation}
		\label{eq:brownian-martingale-L2}
		\mathbb E\int_0^T
		|\nabla u^n(s,Y_{s-})\sigma_s|^2\,ds
		\leq
		L_f^2K_\sigma^2T
		<\infty.
	\end{equation}
	Since \(u^n(t,\cdot)\) is globally \(L_f\)-Lipschitz,
	\begin{equation}
		\label{eq:poisson-martingale-L2}
		\begin{aligned}
			&\mathbb E\int_0^T\int_U
			\bigl|
			u^n(s,Y_{s-}+g_s(z))
			-u^n(s,Y_{s-})
			\bigr|^2
			\,\nu(dz)\,ds
			\\
			&\qquad
			\leq
			L_f^2K_g^2T
			<\infty.
		\end{aligned}
	\end{equation}
	Thus \(M^{W,n}\) and \(M^{N,n}\) are square-integrable martingales,
	not merely local martingales. Since \(\gamma\) and \(\tau\) are
	bounded, optional sampling yields
	\begin{equation}
		\label{eq:conditional-martingale-increments-zero}
		\mathbb E\left[
		M_\tau^{W,n}-M_\gamma^{W,n}
		+
		M_\tau^{N,n}-M_\gamma^{N,n}
		\,\middle|\,
		\mathcal F_\gamma
		\right]
		=0.
	\end{equation}

	

	\textbf{Step 3: It\^o--L\'evy formula and the smooth estimate.}
	
	To justify the It\^o--L\'evy formula without assuming a global bound
	on \(D^2u^n\), choose
	\[
	\chi\in C_c^\infty(\mathbb R^d),
	\qquad
	0\leq\chi\leq1,
	\qquad
	\chi=1\ \text{on }B_1,
	\qquad
	\operatorname{supp}\chi\subset B_2,
	\]
	and set
	\[
	\chi_m(x):=\chi(x/m),
	\qquad
	u_m^n(t,x):=\chi_m(x)u^n(t,x).
	\]
	Then
	\[
	\nabla u_m^n
	=
	\chi_m\nabla u^n
	+
	u^n\nabla\chi_m.
	\]
	Since
	\[
	\nabla\chi_m(x)
	=
	\frac1m\nabla\chi(x/m)
	\]
	and \(\nabla\chi_m\) is supported in
	\(\{m\leq|x|\leq2m\}\), the growth estimate
	\eqref{eq:regularized-growth-gradient-bounds} gives
	\[
	\begin{aligned}
		|u^n(t,x)|\,|\nabla\chi_m(x)|
		&\leq
		L_f(1+2m)
		\frac{\|\nabla\chi\|_\infty}{m}
		\\
		&\leq
		C_\chi L_f
	\end{aligned}
	\]
	on the support of \(\nabla\chi_m\). Hence
	\begin{equation}
		\label{eq:truncated-uniform-gradient}
		\sup_{m\geq1}
		\|\nabla u_m^n\|_{L^\infty(\mathbb R\times\mathbb R^d)}
		\leq
		C_\chi L_f.
	\end{equation}
	
	For
	\[
	r_{m,n}(t,x,h)
	:=
	u_m^n(t,x+h)
	-u_m^n(t,x)
	-\nabla u_m^n(t,x)\cdot h,
	\]
	\eqref{eq:truncated-uniform-gradient} implies
	\[
	|r_{m,n}(t,x,h)|
	\leq
	C_\chi L_f|h|
	\leq
	C_\chi L_f|h|^2,
	\qquad |h|>1.
	\]
	If \(m>R+1\), \(x\in B_R\), and \(|h|\leq1\), then
	\(x,x+h\in B_{R+1}\subset B_m\), and therefore
	\[
	u_m^n(t,x)=u^n(t,x),
	\qquad
	u_m^n(t,x+h)=u^n(t,x+h),
	\qquad
	\nabla u_m^n(t,x)=\nabla u^n(t,x).
	\]
	Consequently,
	\[
	0
	\leq
	r_{m,n}(t,x,h)
	=
	r_n(t,x,h)
	\leq
	\frac12H_{n,R}|h|^2.
	\]
	Thus, for \(m>R+1\) and \(x\in B_R\),
	\begin{equation}
		\label{eq:truncated-jump-dominator}
		|r_{m,n}(t,x,h)|
		\leq
		\frac12H_{n,R}|h|^2
		\mathbf1_{\{|h|\leq1\}}
		+
		C_\chi L_f|h|^2
		\mathbf1_{\{|h|>1\}}.
	\end{equation}
	After substituting \(h=g_s(z)\), the right-hand side is integrable
	with respect to
	\(\nu(dz)\,ds\,d\mathbb P\) by
	\eqref{eq:ito-levy-uniform-bounds}.
	
Apply the It\^o--L\'evy formula to \(u_m^n(t,Y_t)\) on
\([\gamma,\tau]\) and take conditional expectation with respect to
\(\mathcal F_\gamma\). By
\eqref{eq:conditional-martingale-increments-zero}, the Brownian and
compensated Poisson martingale increments have conditional expectation
zero. It remains to pass to the limit \(m\to\infty\) in the remaining
terms.

Since \(\tau\leq\tau_R(Y)\), we have
\[
Y_{s-}\in\overline B_R,
\qquad
\gamma<s<\tau.
\]
If \(m>R+1\), then \(\chi_m\equiv1\) on a neighborhood of
\(\overline B_R\). Hence, on \(\{\gamma<s<\tau\}\),
\[
\partial_su_m^n(s,Y_{s-})
=
\partial_su^n(s,Y_{s-}),
\]
and likewise
\[
\nabla u_m^n(s,Y_{s-})
=
\nabla u^n(s,Y_{s-}),
\qquad
D^2u_m^n(s,Y_{s-})
=
D^2u^n(s,Y_{s-}).
\]
Thus the time-derivative, drift, and continuous second-order terms
already coincide with those for \(u^n\) once \(m>R+1\).

For the jump compensator, for every fixed \((s,\omega,z)\),
\[
r_{m,n}\bigl(s,Y_{s-},g_s(z)\bigr)
\longrightarrow
r_n\bigl(s,Y_{s-},g_s(z)\bigr)
\qquad\text{as }m\to\infty.
\]
Together with the uniform bound
\eqref{eq:truncated-jump-dominator}, dominated convergence yields
\[
\mathcal J_su_m^n(s,Y_{s-})
\longrightarrow
\mathcal J_su^n(s,Y_{s-})
\]
in the corresponding integrable sense, and hence also after integration
over \([\gamma,\tau]\) and conditional expectation.

Finally,
\[
|u_m^n(t,x)|
\leq
|u^n(t,x)|
\leq
L_f(1+|x|)
\]
for all \(m\). Since the boundedness assumptions imply
\(Y_\gamma,Y_\tau\in L^1\), while
\[
u_m^n(\gamma,Y_\gamma)\to u^n(\gamma,Y_\gamma),
\qquad
u_m^n(\tau,Y_\tau)\to u^n(\tau,Y_\tau)
\]
almost surely, conditional dominated convergence applies to the endpoint
terms.

Letting \(m\to\infty\) therefore gives
\begin{equation}
	\label{eq:ito-formula-krylov-revised}
	\begin{aligned}
		&\mathbb E\left[
		u^n(\tau,Y_\tau)
		-u^n(\gamma,Y_\gamma)
		\,\middle|\,
		\mathcal F_\gamma
		\right]
		\\
		=&
		\mathbb E\left[
		\left.
		\int_\gamma^\tau
		\Bigl[
		\partial_su^n(s,Y_{s-})
		+b_s\cdot\nabla u^n(s,Y_{s-})
		\right.\right.
		\\
		&\left.\left.
		+\operatorname{tr}\!\left(
		A_sD^2u^n(s,Y_{s-})
		\right)
		+\mathcal J_su^n(s,Y_{s-})
		\Bigr]\,ds
		\,\right|\,
		\mathcal F_\gamma
		\right].
	\end{aligned}
\end{equation}

	For \(ds\,d\mathbb P\)-almost every \((s,\omega)\) with
	\(\gamma<s<\tau\), one has \(Y_{s-}\in B_R\). Therefore, by
	\eqref{eq:regularized-aleksandrov-inequality},
	\eqref{eq:jump-operator-positive}, and \(|b_s|\leq K_b\),
	\begin{equation}
		\label{eq:generator-lower-bound-revised}
		\begin{aligned}
			&\partial_su^n(s,Y_{s-})
			+b_s\cdot\nabla u^n(s,Y_{s-})
			+\operatorname{tr}\!\left(
			A_sD^2u^n(s,Y_{s-})
			\right)
			+\mathcal J_su^n(s,Y_{s-})
			\\
			\geq&
			\kappa_d(\det A_s)^{1/(d+1)}
			f^n(s,Y_{s-})
			-K_bL_f.
		\end{aligned}
	\end{equation}
	Since a c\`adl\`ag path has at most countably many jumps,
	\(Y_s=Y_{s-}\) for Lebesgue-almost every \(s\). Combining
	\eqref{eq:ito-formula-krylov-revised} and
	\eqref{eq:generator-lower-bound-revised} gives
	\begin{equation}
		\label{eq:preliminary-weighted-bound-revised}
		\begin{aligned}
			&\kappa_d
			\mathbb E\left[
			\left.
			\int_\gamma^\tau
			(\det A_s)^{1/(d+1)}
			f^n(s,Y_s)\,ds
			\,\right|\,
			\mathcal F_\gamma
			\right]
			\\
			\leq&
			\mathbb E\left[
			u^n(\tau,Y_\tau)
			-u^n(\gamma,Y_\gamma)
			\,\middle|\,
			\mathcal F_\gamma
			\right]
			+K_bTL_f.
		\end{aligned}
	\end{equation}
	
	Let
	\[
	E:=\{\gamma<\tau\}\in\mathcal F_\gamma.
	\]
	On \(E\), one has \(\gamma<\tau_R(Y)\), and hence
	\(Y_\gamma\in B_R\), so \(|Y_\gamma|\leq R\).
	Moreover, conditional Cauchy--Schwarz and the isometries for the
	Brownian and compensated Poisson integrals imply
	\begin{equation}
		\label{eq:conditional-terminal-moment-revised}
		\begin{aligned}
			\mathbf1_E
			\mathbb E\left[
			|Y_\tau|
			\,\middle|\,
			\mathcal F_\gamma
			\right]
			&\leq
			\mathbf1_E
			\left(
			|Y_\gamma|
			+K_bT
			+K_\sigma T^{1/2}
			+K_gT^{1/2}
			\right)
			\\
			&\leq
			\mathbf1_E
			\left[
			R+K_bT
			+(K_\sigma+K_g)T^{1/2}
			\right].
		\end{aligned}
	\end{equation}
	In particular,
	\eqref{eq:regularized-growth-gradient-bounds} yields
	\begin{equation}
		\label{eq:terminal-test-function-bound-revised}
		\mathbf1_E
		\left|
		\mathbb E\left[
		u^n(\tau,Y_\tau)
		-u^n(\gamma,Y_\gamma)
		\,\middle|\,
		\mathcal F_\gamma
		\right]
		\right|
		\leq
		C\mathbf1_E
		\|f\|_{L^{d+1}(Q_{T,R})},
	\end{equation}
	where
	\[
	C=C(d,T,R,K_b,K_\sigma,K_g)
	\]
	is independent of \(n\), \(\gamma\), \(\tau\), and \(f\).
	
	On \(E^c=\{\gamma=\tau\}\), the occupation integral is zero.
	Hence \eqref{eq:preliminary-weighted-bound-revised} and
	\eqref{eq:terminal-test-function-bound-revised} imply
	\begin{equation}
		\label{eq:weighted-krylov-regularized}
		\mathbb E\left[
		\left.
		\int_\gamma^\tau
		(\det A_s)^{1/(d+1)}
		f^n(s,Y_s)\,ds
		\,\right|\,
		\mathcal F_\gamma
		\right]
		\leq
		C_{\mathrm K}
		\|f\|_{L^{d+1}(Q_{T,R})}.
	\end{equation}
	The upper bound on \(\sigma\) gives
	\begin{equation}
		\label{eq:determinant-upper-bound}
		(\det A_s)^{1/(d+1)}
		\leq
		C(d,K_\sigma).
	\end{equation}
	Therefore \eqref{eq:regularized-f-convergence} permits passage to the
	limit in \eqref{eq:weighted-krylov-regularized}, and we obtain
	\begin{equation}
		\label{eq:weighted-krylov-smooth-revised}
		\mathbb E\left[
		\left.
		\int_\gamma^\tau
		(\det A_s)^{1/(d+1)}
		f(s,Y_s)\,ds
		\,\right|\,
		\mathcal F_\gamma
		\right]
		\leq
		C_{\mathrm K}
		\|f\|_{L^{d+1}(Q_{T,R})}
	\end{equation}
	for every non-negative
	\(f\in C_c^\infty((0,T)\times B_R)\).
	
	
	\textbf{Step 4: extension from smooth functions to	\(L^{d+1}(Q_{T,R})\).}
	
	For \(A\in\mathcal F_\gamma\), define a finite Borel measure	\(\Gamma_A\) on \(Q_{T,R}\) by
	\begin{equation}
		\label{eq:conditional-occupation-measure}
		\Gamma_A(B)	:=
		\mathbb E\left[
		\mathbf1_A
		\int_\gamma^\tau
		(\det A_s)^{1/(d+1)}
		\mathbf1_B(s,Y_s)\,ds
		\right],
		\qquad
		B\in\mathcal B(Q_{T,R}).
	\end{equation}
	Finiteness follows from
	\eqref{eq:determinant-upper-bound}. Moreover,
	\[
	\Gamma_A\bigl(\{0,T\}\times B_R\bigr)=0.
	\]
	Set
	\[
	D_{T,R}:=(0,T)\times B_R.
	\]
	Thus \(\Gamma_A\) is determined by its restriction to \(D_{T,R}\),
	and
	\[
	L^{d+1}(D_{T,R})
	=
	L^{d+1}(Q_{T,R})
	\]
	up to the usual identification of functions on Lebesgue-null sets.
	
	Multiplying \eqref{eq:weighted-krylov-smooth-revised} by
	\(\mathbf1_A\) and taking expectations gives
	\begin{equation}
		\label{eq:occupation-measure-smooth-bound}
		\int_{D_{T,R}}\varphi\,d\Gamma_A
		\leq
		C_{\mathrm K}\mathbb P(A)
		\|\varphi\|_{L^{d+1}(D_{T,R})}
	\end{equation}
	for every non-negative
	\(\varphi\in C_c^\infty(D_{T,R})\).
	
	Since \(D_{T,R}\) is open, every non-negative
	\(\varphi\in C_c(D_{T,R})\) can be uniformly approximated by
	non-negative functions in \(C_c^\infty(D_{T,R})\).
	Consequently, \eqref{eq:occupation-measure-smooth-bound} also holds
	for every non-negative \(\varphi\in C_c(D_{T,R})\).
	Applying this estimate to \(\varphi^+\) and \(\varphi^-\) gives
	\[
	\left|
	\int_{D_{T,R}}\varphi\,d\Gamma_A
	\right|
	\leq
	C_{\mathrm K}\mathbb P(A)
	\|\varphi\|_{L^{d+1}(D_{T,R})},
	\qquad
	\varphi\in C_c(D_{T,R}).
	\]
	
	Since \(C_c^\infty(D_{T,R})\) is dense in
	\(L^{d+1}(D_{T,R})\), the functional
	\[
	\varphi
	\longmapsto
	\int_{D_{T,R}}\varphi\,d\Gamma_A
	\]
	extends uniquely to a bounded linear functional on
	\(L^{d+1}(D_{T,R})\). By
	\(L^p\)--\(L^{p'}\) duality, there exists
	\[
	h_A
	\in
	L^{(d+1)/d}(D_{T,R})
	\]
	such that
	\[
	\int_{D_{T,R}}\varphi\,d\Gamma_A
	=
	\int_{D_{T,R}}
	\varphi(t,x)h_A(t,x)\,dt\,dx
	\]
	for \(\varphi\in C_c(D_{T,R})\), and
	\[
	\|h_A\|_{L^{(d+1)/d}(D_{T,R})}
	\leq
	C_{\mathrm K}\mathbb P(A).
	\]
	The finite Radon measures
	\(\Gamma_A|_{D_{T,R}}\) and
	\(h_A(t,x)\,dt\,dx\) agree on \(C_c(D_{T,R})\).
	By uniqueness in the Riesz representation theorem,
	\[
	\Gamma_A|_{D_{T,R}}(dt,dx)
	=
	h_A(t,x)\,dt\,dx.
	\]
	Since \(\Gamma_A\) gives no mass to
	\(\{0,T\}\times B_R\), we may equivalently write, after extending
	\(h_A\) arbitrarily on this Lebesgue-null set,
	\begin{equation}
		\label{eq:occupation-density}
		\Gamma_A(dt,dx)
		=
		h_A(t,x)\,dt\,dx
		\qquad\text{on }Q_{T,R}.
	\end{equation}
	
	Therefore, for every non-negative Borel function
	\(f\in L^{d+1}(Q_{T,R})\), H\"older's inequality gives
	\[
	\begin{aligned}
		\int_{Q_{T,R}}f\,d\Gamma_A
		&=
		\int_{Q_{T,R}}
		f(t,x)h_A(t,x)\,dt\,dx
		\\
		&\leq
		\|f\|_{L^{d+1}(Q_{T,R})}
		\|h_A\|_{L^{(d+1)/d}(Q_{T,R})}
		\\
		&\leq
		C_{\mathrm K}\mathbb P(A)
		\|f\|_{L^{d+1}(Q_{T,R})}.
	\end{aligned}
	\]
	By the definition of \(\Gamma_A\), first for indicator functions
	and then for non-negative Borel functions by monotone convergence,
	this is precisely
	\begin{equation}
		\label{eq:occupation-measure-general-bound}
		\mathbb E\left[
		\mathbf1_A
		\int_\gamma^\tau
		(\det A_s)^{1/(d+1)}
		f(s,Y_s)\,ds
		\right]
		\leq
		C_{\mathrm K}\mathbb P(A)
		\|f\|_{L^{d+1}(Q_{T,R})}.
	\end{equation}
	
	Since \eqref{eq:occupation-measure-general-bound} holds for every
	\(A\in\mathcal F_\gamma\), the defining property of conditional
	expectation gives
	\eqref{eq:weighted-conditional-krylov}. Indeed, if
	\[
	\mathbb E\left[
	\left.
	\int_\gamma^\tau
	(\det A_s)^{1/(d+1)}
	f(s,Y_s)\,ds
	\,\right|\,
	\mathcal F_\gamma
	\right]
	>
	C_{\mathrm K}
	\|f\|_{L^{d+1}(Q_{T,R})}
	\]
	on some set \(A\in\mathcal F_\gamma\) of positive probability, then
	integrating over \(A\) would contradict
	\eqref{eq:occupation-measure-general-bound}.
\end{proof}

\end{document}
